% =====================================================================
% Capítulo 3 — Experimento preliminar: mariposas
% =====================================================================
\chapter{Generando mariposas: modelos de difusión en acción}
\label{cap:mariposas}

Como primer experimento práctico, se entrena un modelo DDPM
(\textit{Denoising Diffusion Probabilistic Model}) sobre un
conjunto de imágenes de mariposas, con el objetivo de generar
muestras sintéticas nuevas. Este ejercicio sirve como banco de
pruebas para comprender el proceso de entrenamiento e inferencia
antes de abordar el problema de reconstrucción de imágenes PET.

\section{Datos y preprocesamiento}
\label{sec:mariposas-datos}

Se utiliza el subconjunto \texttt{huggan/smithsonian\_butterflies\_subset}
del Instituto Smithsoniano. Cada imagen se redimensiona a
$256 \times 256$ píxeles preservando su relación de aspecto: el
lado más largo se reescala al tamaño objetivo, y el lado más corto
se rellena con relleno blanco simétrico hasta alcanzar una imagen
cuadrada (operación \texttt{ResizeToFit} + \texttt{PadToSquare} en
el código fuente). Sobre el resultado se aplica volteo horizontal
aleatorio como aumento de datos, y se normaliza al rango $[-1, 1]$
mediante la transformación
\begin{equation}
    x' = \frac{x - 0.5}{0.5}.
\end{equation}
El conjunto se carga en lotes de orden aleatorio en cada época.

Una iteración previa de este experimento siguió literalmente el
\textit{tutorial} oficial de la librería \texttt{diffusers}, con
imágenes reescaladas directamente a $128 \times 128$ píxeles. Ese
reescalado directo (sin preservar relación de aspecto ni rellenar
para mantener la geometría original) introduce una deformación
sistemática del contenido, visible en las muestras generadas
(\Cref{fig:mariposas-128}). El paso a $256 \times 256$ con
\texttt{ResizeToFit} + \texttt{PadToSquare} corrige esta patología
y es el preprocesamiento empleado en el resto del capítulo.

\section{Paradigma de difusión empleado}
\label{sec:mariposas-paradigma}

El experimento sigue la formulación DDPM original en su versión más
sencilla, que es exactamente la desarrollada en el
\Cref{cap:marco-teorico}:

\begin{itemize}
    \item \textbf{Proceso hacia delante}: corrupción gaussiana en
    $T = 1000$ pasos discretos con calendario de varianza
    \emph{lineal}, muestreando $\mathbf{x}_t$ directamente desde
    $\mathbf{x}_0$ mediante la forma cerrada de las
    \Cref{eq:forward-closed,eq:forward-reparam}.
    \item \textbf{Objetivo de entrenamiento}: predicción del ruido
    ($\epsilon$-prediction) con la pérdida MSE simplificada de la
    \Cref{eq:loss-eps}. En cada iteración se muestrea $t$
    uniformemente en $\{1, \ldots, 1000\}$ y el gradiente se propaga
    únicamente a través de la predicción del modelo.
    \item \textbf{Inferencia}: muestreo ancestral DDPM completo de
    $T = 1000$ pasos según la \Cref{eq:ddpm-step}, partiendo de
    ruido puro $\mathbf{x}_T \sim \mathcal{N}(\mathbf{0},
    \mathbf{I})$. La implementación permite fijar la semilla para
    reproducibilidad o variarla para explorar el espacio generativo.
\end{itemize}

Esta combinación —$\epsilon$-prediction, calendario lineal, muestreo
ancestral— constituye la línea base canónica del paradigma. El
\Cref{cap:pet} parte de ella y la mejora con las elecciones
discutidas en la \Cref{subsec:parametrizaciones} (parametrización
$\mathbf{v}$, calendario cosenoidal, muestreo DDIM).

\section{Arquitectura del modelo}
\label{sec:mariposas-arquitectura}

La red neuronal $\boldsymbol{\varepsilon}_\theta(\mathbf{x}_t, t)$
encargada de predecir el ruido es una \emph{U-Net}, la arquitectura
convolucional que se ha consolidado como columna vertebral de
prácticamente todos los modelos de difusión sobre imágenes. Conviene
detenerse brevemente en por qué esta arquitectura —y no otra— es la
elección natural para la tarea de \textit{denoising}.

La U-Net, introducida por Ronneberger et
al.~\cite{ronneberger2015unet} para segmentación de imágenes
biomédicas, es una red \emph{codificador--decodificador} con forma de
«U» (\Cref{fig:unet}). El \textbf{camino de contracción} (rama
izquierda) alterna convoluciones y submuestreos que reducen
progresivamente la resolución espacial mientras aumentan el número de
canales: la red comprime la imagen en una representación cada vez más
abstracta y de menor resolución, capturando contexto a gran escala.
El \textbf{camino de expansión} (rama derecha) hace lo inverso,
recuperando la resolución mediante sobremuestreos hasta producir una
salida con las mismas dimensiones espaciales que la entrada.

El elemento distintivo son las \textbf{conexiones de salto}
(\textit{skip connections}, flechas grises horizontales en la
\Cref{fig:unet}): cada nivel del decodificador recibe, concatenado,
el mapa de características del nivel homólogo del codificador. Estas
conexiones reinyectan el detalle espacial de alta frecuencia —bordes,
texturas, localización precisa— que de otro modo se perdería en el
cuello de botella. Es justamente esta propiedad la que hace a la
U-Net idónea para los problemas de predicción densa píxel a píxel:
la salida debe tener la misma resolución que la entrada y preservar
su estructura fina, combinando contexto global (\emph{qué} hay en la
imagen) con detalle local (\emph{dónde} está exactamente). En nuestro
caso esa salida no es un mapa de segmentación, sino la predicción del
ruido $\boldsymbol{\varepsilon}$ —y, en el experimento PET, de
$\mathbf{v}$—, una tarea que comparte exactamente la misma exigencia
estructural.

\begin{figure}[h!]
    \centering
    \includegraphics[width=0.9\linewidth]{img/U-net_example.png}
    \caption{Arquitectura U-Net original de Ronneberger et
    al.~\cite{ronneberger2015unet}. El camino de contracción
    (izquierda) reduce la resolución y aumenta los canales; el de
    expansión (derecha) la restaura mediante \textit{up-convolutions};
    las flechas grises horizontales son las conexiones de salto que
    reinyectan el detalle espacial del codificador en el decodificador.
    El ejemplo corresponde a la tarea de segmentación para la que se
    concibió (salida de dos canales de clase); las U-Net de difusión
    conservan esta topología en «U» pero añaden \textit{embeddings} de
    paso temporal y bloques de atención, y producen una salida de la
    misma naturaleza que la entrada.}
    \label{fig:unet}
\end{figure}

Frente a la U-Net original —concebida para segmentación, con una
salida de pocos canales de clase y recortes de borde
(\textit{valid padding})—, las U-Net de difusión incorporan dos
adaptaciones clave: \textbf{embeddings de paso temporal}, que informan
a la red del nivel de ruido en el que opera, y \textbf{bloques
residuales y de atención} en lugar de las convoluciones simples del
diseño de 2015.

Nuestra implementación concreta (un \texttt{UNet2DModel} de la
librería \texttt{diffusers}) presenta seis niveles de
resolución con canales $(128, 128, 256, 256, 512, 512)$ y dos
capas ResNet por bloque. La atención multi-cabeza se incorpora
únicamente en el quinto nivel, tanto en el camino descendente
(\texttt{AttnDownBlock2D}) como en el ascendente
(\texttt{AttnUpBlock2D}), donde la resolución espacial es
suficientemente reducida para que el coste cuadrático de la
atención sea asumible. El paso temporal $t$ se inyecta en cada
bloque ResNet mediante \textit{embeddings} sinusoidales, análogos a
los empleados en el Transformer original, de modo que el modelo
aprende una política de \textit{denoising} condicional en el nivel
de ruido.

\section{Configuración del entrenamiento}
\label{sec:mariposas-config}

El modelo se optimiza con AdamW con tasa de aprendizaje inicial
$\eta = 10^{-4}$ y un planificador cosenoidal con 500 pasos de
calentamiento (\textit{warmup}). El entrenamiento se lleva a cabo
durante 100 épocas con precisión mixta \texttt{fp16} para reducir
el consumo de memoria, y se recortan los gradientes a norma
máxima 1.0 para garantizar la estabilidad numérica. Toda esta
fase preliminar se ejecuta \textbf{localmente} sobre un MacBook
Pro M4 Pro con aceleración MPS, suficiente para entrenar el
modelo a $256 \times 256$ en un plazo razonable y validar el
flujo de entrenamiento e inferencia antes de pasar a la tarea de
reconstrucción PET (véase \Cref{sec:pet-entorno-cloud}).

\section{Aprendizajes y resultados}
\label{sec:mariposas-aprendizajes}

Los modelos de difusión son buenos generalizando, ya que es más
complicado que puedan memorizar los ejemplos de entrenamiento,
dado que van a ver las mismas imágenes pero con una cantidad de
ruido completamente diferente entre una época y otra. Se ven
forzados a aprender la distribución subyacente de las imágenes
que analizan, en este caso aprendiendo cómo es una mariposa.
Aprender a predecir la cantidad de ruido que hay en una imagen
es matemáticamente equivalente a predecir la imagen de una
mariposa a partir de una imagen ruidosa.

Se ha observado que la pérdida durante el entrenamiento decrece
muy rápidamente en las primeras épocas, y después se estabiliza
a lo largo de todo el proceso de entrenamiento. Esto no significa
que el modelo deje de aprender una vez se estabiliza la pérdida,
sino que aprende a predecir la cantidad de ruido que hay en las
imágenes más rápidamente; pero los parámetros se siguen
actualizando y aprendiendo con más certeza la distribución
subyacente de imágenes de mariposas. Esta observación —que la
pérdida de entrenamiento es una mala señal de convergencia para
los modelos de difusión— se convertirá en una decisión de diseño
en el experimento principal
(\Cref{subsec:pet-config}).

Para ilustrar este fenómeno comparamos las muestras generadas por
el modelo en dos momentos opuestos del entrenamiento: las primeras
20 épocas y las últimas 20 épocas (cada panel corresponde al
\textit{checkpoint} guardado al final de la época indicada).

La \Cref{fig:mariposas-tempranas} recoge las muestras de las
primeras 20 épocas, concretamente las correspondientes a los
\textit{checkpoints} de las épocas~10 y~20. En esta fase temprana
las mariposas apenas son reconocibles: aparecen como manchas de
color difusas, sin contornos definidos ni simetría, y con
texturas dominadas por artefactos de alta frecuencia. El modelo ya
ha captado la paleta cromática general del conjunto de datos
(tonos cálidos sobre fondos vegetales), pero todavía no ha
aprendido la estructura espacial que caracteriza a una mariposa.

\begin{figure}[h!]
    \centering
    \includegraphics[width=0.35\linewidth]{img/generated_butterflies/256pix_0009.png}\hspace{1em}
    \includegraphics[width=0.35\linewidth]{img/generated_butterflies/256pix_0019.png}
    \caption{Muestras de las primeras 20 épocas a resolución
    $256 \times 256$: \textit{checkpoints} al final de las
    épocas~10 (izquierda) y~20 (derecha). Las muestras son borrosas
    y carecen de estructura coherente, aunque ya reproducen la
    paleta cromática del conjunto de entrenamiento.}
    \label{fig:mariposas-tempranas}
\end{figure}

En contraste, la \Cref{fig:mariposas-tardias} muestra las muestras
de las últimas 20 épocas, correspondientes a los
\textit{checkpoints} de las épocas~90 y~100. A pesar de que la
pérdida es prácticamente la misma que en épocas más tempranas, la
calidad perceptual ha mejorado de forma notable: la red genera
mariposas con siluetas reconocibles, simetría bilateral aproximada
y patrones de color coherentes con la distribución de
entrenamiento. Se observa que las imágenes empiezan a ser buenas
en torno a la época~90 y se estabilizan en la época~100.

\begin{figure}[h!]
    \centering
    \includegraphics[width=0.35\linewidth]{img/generated_butterflies/256pix_0089.png}\hspace{1em}
    \includegraphics[width=0.35\linewidth]{img/generated_butterflies/256pix_0099.png}
    \caption{Muestras de las últimas 20 épocas a resolución
    $256 \times 256$: \textit{checkpoints} al final de las
    épocas~90 (izquierda) y~100 (derecha). Las muestras exhiben
    siluetas y simetrías reconocibles y patrones de color
    consistentes con la distribución de entrenamiento.}
    \label{fig:mariposas-tardias}
\end{figure}

\paragraph{Comparación en inferencia: $128 \times 128$ frente a
$256 \times 256$.}
Una vez analizada la evolución del entrenamiento, comparamos la
calidad de las muestras obtenidas en inferencia con las dos
resoluciones consideradas. La \Cref{fig:mariposas-128} muestra dos
ejemplos generados por la iteración previa del experimento,
entrenada con reescalado directo a $128 \times 128$ sin preservar
la relación de aspecto. Aunque la red aprende ciertas
regularidades de color y forma, las mariposas aparecen
visiblemente deformadas, lo cual ilustra el impacto del
preprocesamiento sobre la calidad final.

\begin{figure}[h!]
    \centering
    \includegraphics[width=0.35\linewidth]{img/generated_butterflies/128pix_sample1.jpeg}\hspace{1em}
    \includegraphics[width=0.35\linewidth]{img/generated_butterflies/128pix_sample2.jpeg}
    \caption{Muestras en inferencia a resolución $128 \times 128$
    (reescalado directo sin \textit{padding}). Las muestras
    presentan deformaciones sistemáticas heredadas del
    preprocesamiento.}
    \label{fig:mariposas-128}
\end{figure}

Por su parte, la \Cref{fig:mariposas-256} recoge dos ejemplos
generados en inferencia a resolución $256 \times 256$. Frente a la
línea base anterior, estas muestras presentan siluetas mucho más
nítidas y proporciones realistas, sin las deformaciones
sistemáticas de la versión a $128 \times 128$. El aumento de
resolución, combinado con un preprocesamiento que preserva la
relación de aspecto, se traduce en mariposas con mayor nivel de
detalle y mejor definición de los patrones alares.

\begin{figure}[h!]
    \centering
    \includegraphics[width=0.35\linewidth]{img/generated_butterflies/256pix_sample1.jpeg}\hspace{1em}
    \includegraphics[width=0.35\linewidth]{img/generated_butterflies/256pix_sample2.jpeg}
    \caption{Muestras en inferencia a resolución $256 \times 256$.
    Frente a la línea base de $128 \times 128$, las muestras son
    más nítidas, con proporciones realistas y mayor nivel de
    detalle en los patrones alares.}
    \label{fig:mariposas-256}
\end{figure}
